\chapter{Vom Prompt zum Arbeitsauftrag}
Ein klassischer Prompt verlangt eine Antwort. Ein moderner Arbeitsauftrag beschreibt dagegen ein Ergebnis, Randbedingungen und Qualitätskriterien. Das System kann daraus Teilschritte ableiten, Werkzeuge wählen, Zwischenergebnisse prüfen und ein nutzbares Artefakt erzeugen.
\begin{center}\begin{tikzpicture}[node distance=15mm]
\node[box] (u) {Nutzerziel}; \node[box,right=of u] (o) {Orchestrierung}; \node[box,right=of o] (e) {Ergebnis};
\draw[arr] (u)--node[above]{Auftrag}(o); \draw[arr] (o)--node[above]{geprüft}(e);
\node[tool,below=of o,xshift=-28mm] (m) {Modell};\node[tool,below=of o] (w) {Werkzeuge};\node[tool,below=of o,xshift=28mm] (d) {Daten};
\draw[arr] (o)--(m);\draw[arr] (o)--(w);\draw[arr] (o)--(d);
\end{tikzpicture}\end{center}

\begin{merke}Ein guter Auftrag nennt Ziel, Kontext, Grenzen und Erfolgskriterien. Er muss nicht jeden Mausklick vorgeben.\end{merke}
\section{Vier Ebenen eines guten Auftrags}
\begin{enumerate}\item \textbf{Ziel:} Was soll am Ende vorliegen?\item \textbf{Kontext:} Für wen und wofür?\item \textbf{Grenzen:} Budget, Zeit, Quellen, Format und Berechtigungen.\item \textbf{Prüfung:} Woran erkennt man ein gutes Ergebnis?\end{enumerate}
\begin{praxis}Vergleiche: ``Erkläre RAG'' mit ``Erstelle für einen Softwareentwickler einen 20-minütigen Lernpfad zu RAG, mit Mini-Beispiel, typischen Fehlern und Selbsttest.''\end{praxis}
